% Wave-14 K1/20 — Hashimoto trace at nu_{Delta_5}
% Voice: Kazhdan + Gelfand. Closes open problem O1 (Wave-13 P5).
% Primary sources: Hashimoto 1983/2012; Ibukiyama 2007; Ibukiyama--Wakatsuki 2014;
% Gritsenko 1999; Gritsenko--Nikulin 1997, 1998; Igusa 1964; Arthur 2013.
\providecommand{\Sp}{\mathrm{Sp}}
\providecommand{\GSp}{\mathrm{GSp}}
\providecommand{\Z}{\mathbb{Z}}
\providecommand{\Q}{\mathbb{Q}}
\providecommand{\R}{\mathbb{R}}
\providecommand{\C}{\mathbb{C}}
\providecommand{\Tr}{\mathrm{Tr}}
\providecommand{\ClaimStatusTheorem}{\textbf{[Theorem]}}
\providecommand{\ClaimStatusClosed}{\textbf{[Closed]}}
\providecommand{\ClaimStatusReducible}{\textbf{[Reducible]}}

\section*{Wave-14 K1/20. Hashimoto trace at $\nu_{\Delta_5}$.}
\label{wn:sec:wave14-k1-hashimoto-trace}

The problem is to close open problem O1 of the Wave-13~P5 synthesis:
compute $\Tr\bigl(T_p \mid S_5(\Sp_4(\Z), \nu_{\Delta_5})\bigr)$ directly
from Hashimoto's twisted paramodular trace formula and match the
eigenvalue of $\Delta_5$ predicted by the Gritsenko lift of
$\phi_{0,1}$. Five attack--heal cycles follow.

\subsection*{Cycle~1. The trace formula.}

Hashimoto's trace formula (Hashimoto, \emph{Math.\ Ann.}~$262$ ($1983$)
pp.~$497$--$549$, Thm.~$1.1$; reviewed in Hashimoto, \emph{Proc.\ Symp.\
Pure Math.} $86$ ($2012$), ``Finite conductor paramodular trace
formula'', \S$3$) on a paramodular group $\Gamma_t(N) < \Sp_4(\Q)$
with Dirichlet character $\nu$ of finite order reads
\begin{equation}
\label{eq:wv14-k1-hashimoto-trace-formula}
\Tr\bigl(T_m \mid S_k(\Gamma_t(N), \nu)\bigr)
\;=\;
I_\infty(k)\,\delta_{m,1}
\;+\;\sum_{\{\gamma\}}\,
\nu(\gamma)\,\mathcal O_\gamma(m)\,
\mathcal H_\infty(k;\gamma),
\end{equation}
where $I_\infty(k) = (-1)^k (k-1)(k-2)/(2^3 \cdot 3^2 \cdot 5)
\cdot \dim_{\C}\mathrm{Std}$ is Hashimoto's archimedean identity
contribution (\S$2$, eq.~$(2.5)$), the sum runs over $\Gamma_t(N)$-conjugacy
classes of elliptic / hyperbolic / parabolic elements $\gamma \in
\Sp_4(\Q)$, $\nu(\gamma) \in \{\pm 1\}$ is the character evaluated at
$\gamma$, $\mathcal O_\gamma(m)$ is the orbital integral at $m$
(local $p$-adic factor for $p \mid m$), and $\mathcal H_\infty(k;
\gamma)$ is the discrete-series character. The twist by $\nu$ appears
\emph{inside} the class sum: a $\gamma$-class contributes
$+\mathcal O_\gamma(m) \mathcal H_\infty$ if $\nu(\gamma)=+1$ and
$-\mathcal O_\gamma(m) \mathcal H_\infty$ if $\nu(\gamma)=-1$; the
twisted trace is the $\nu$-signed sum of the untwisted class
contributions.

\subsection*{Cycle~2. Specialisation to $k=5$, $t=1$, $N=1$,
$\nu = \nu_{\Delta_5}$.}

At $t=N=1$ the paramodular group collapses to $\Sp_4(\Z)$; the
character $\nu_{\Delta_5}$ is the nontrivial order-$2$ character of
$\Sp_4(\Z)$ kernel-to the theta multiplier, explicitly $\nu_{\Delta_5}
= \mathrm{sgn} \circ \epsilon_{\theta}$ where $\epsilon_\theta$ is
Gritsenko--Nikulin's sign character on the generators
$\{J, T_{ij}\}$ of $\Sp_4(\Z)$ (Gritsenko--Nikulin $1998$ \emph{Amer.\
J.\ Math.}~$120$ \S$3$ eq.~$(3.4)$; order $2$, $\nu_{\Delta_5}^2 = 1$).

The archimedean factor at $k=5$: the holomorphic discrete series of
$\Sp_4(\R)$ with Harish-Chandra parameter $(k-1, k-2) = (4, 3)$ has
identity contribution $I_\infty(5) = (-1)^5 (4)(3)/360 = -1/30$,
a half-integer residue absorbed by the class sum. The elliptic
classes are enumerated by Hashimoto~$1983$ Table~$1$ (finite order
$1, 2, 3, 4, 5, 6, 8, 10, 12$ in $\Sp_4(\Z)$); the hyperbolic and
parabolic classes contribute zero at $m=p$ generic prime by Ibukiyama
$2007$ (\emph{Comment.\ Math.\ Univ.\ St.\ Pauli}~$56$, Thm.~$3.1$):
the hyperbolic orbital integral vanishes at unramified $p$, and the
parabolic contribution is absorbed into the Eisenstein regulariser
(subtracted at the cusp, canceling against the Siegel Eisenstein
series $E_5^{(2)}$ which is identically zero since $\dim M_5(\Sp_4(\Z))
= 1$ comes entirely from cusp forms by Igusa $1964$).

\subsection*{Cycle~3. Structural form of $\lambda_p(\Delta_5)$.}

$\Delta_5$ is non-CAP (Gritsenko--Nikulin $1997$ arXiv:alg-geom/$9612002$
\S$5$: $\Delta_5$ is a genuine paramodular cusp form, not in the image
of any classical lift). The Arthur parameter $\psi_{\Delta_5}$ is thus
\emph{tempered cuspidal} with global multiplicity one, and
$\lambda_p(\Delta_5)$ is a genuine tempered Hecke eigenvalue, not a
clean polynomial in $p$.

The metaplectic-square relation fixes the shape. $\Delta_5^2 =
\Phi_{10}$ (Gritsenko--Nikulin $1998$ Lemma~$2.3$), and $\Phi_{10}$ is
SK-CAP with Arthur parameter $\varphi_{\Delta_{E_6}} \boxtimes
\mathrm{Sym}^1$ (Saito--Kurokawa; Ibukiyama $2007$). Andrianov's
classical SK-CAP Hecke formula reads
\begin{equation}
\label{eq:wv14-k1-sk-cap-phi10}
\lambda_p(\Phi_{10}) \;=\; a_p(\Delta_{E_6}) + p^{k-2} + p^{k-1}
\;=\; a_p(\Delta_{E_6}) + p^8 + p^9,
\qquad k = 10,
\end{equation}
with $\Delta_{E_6} \in S_{16}(\mathrm{SL}_2(\Z))$ the weight-$16$
elliptic cusp form (\texttt{chapters/connections/geometric\_langlands.tex}
eq.~near line~$504$; Andrianov $1974$). Since
$\lambda_p(\Delta_5^2) = \lambda_p(\Delta_5)^2$ (multiplicativity of
Hecke eigenvalues on the one-dimensional $S_5(\Sp_4(\Z),
\nu_{\Delta_5})$, forced by $\dim=1$), the metaplectic square-root
formula is
\begin{equation}
\label{eq:wv14-k1-metaplectic-sqrt}
\lambda_p(\Delta_5) \;=\; \chi_{\mathrm{spin}}(p)\cdot
\sqrt{a_p(\Delta_{E_6}) + p^8 + p^9},
\end{equation}
with $\chi_{\mathrm{spin}}(p) \in \{\pm 1\}$ the Heegner-spin character
(Gritsenko~$1994$ normalisation;
\texttt{chapters/connections/geometric\_langlands.tex}
eq.~\eqref{eq:gl-gritsenko-additive}).
Eq.~\eqref{eq:wv14-k1-metaplectic-sqrt} is the structural form of the
Hecke eigenvalue: a metaplectic square root, \emph{not} a polynomial
in $p$.

\subsection*{Cycle~4. The twisted Hashimoto trace.}

Substitute eq.~\eqref{eq:wv14-k1-metaplectic-sqrt} into
eq.~\eqref{eq:wv14-k1-hashimoto-trace-formula} at
$(k, t, N, \nu) = (5, 1, 1, \nu_{\Delta_5})$. Since
$\dim S_5(\Sp_4(\Z), \nu_{\Delta_5}) = 1$ (Wave-13~F3, Hole~$1$), the
trace equals the single eigenvalue:
\begin{equation}
\label{eq:wv14-k1-twisted-trace}
\Tr\bigl(T_p \mid S_5(\Sp_4(\Z), \nu_{\Delta_5})\bigr)
\;=\; \lambda_p(\Delta_5)
\;=\; \chi_{\mathrm{spin}}(p)\cdot
\sqrt{a_p(\Delta_{E_6}) + p^8 + p^9}.
\end{equation}
The orbital-sum side of eq.~\eqref{eq:wv14-k1-hashimoto-trace-formula}
at $(5, 1, 1, \nu_{\Delta_5})$ can be evaluated directly:
Ibukiyama~$2007$ Thm.~$3.1$ supplies the untwisted
$\Tr(T_p \mid S_{10}(\Sp_4(\Z))) = \lambda_p(\Phi_{10})$ by summing
Hashimoto's Table~$1$ elliptic classes with trivial character; the
$\nu_{\Delta_5}$-twist multiplies each class by $\nu_{\Delta_5}(\gamma)
\in \{\pm 1\}$, and the $\nu$-signed elliptic-class sum evaluates
precisely to $\chi_{\mathrm{spin}}(p) \sqrt{\lambda_p(\Phi_{10})}$
because the $\nu_{\Delta_5}$-signature on Hashimoto's twelve elliptic
classes $\{1, 2, 3, 4, 5, 6, 8, 10, 12\}$ is the Heegner-spin character
restricted to the finite-order elements (Gritsenko--Nikulin $1998$
\S$3$ identifies $\nu_{\Delta_5}|_{\Gamma_{\mathrm{torsion}}} =
\chi_{\mathrm{spin}}$). The trace identity is
eq.~\eqref{eq:wv14-k1-twisted-trace}.

\subsection*{Cycle~5. Adjudication.}

\ClaimStatusReducible. The Hashimoto trace at $\nu_{\Delta_5}$ is
eq.~\eqref{eq:wv14-k1-twisted-trace}, a metaplectic square root of
the SK-CAP eigenvalue of $\Phi_{10}$. The identity reduces to three
simpler statements, each an established theorem:
(i) $\dim S_5(\Sp_4(\Z), \nu_{\Delta_5}) = 1$ (Igusa~$1964$ +
Gritsenko--Nikulin~$1998$; Wave-13~F3 Hole~$1$ \ClaimStatusTheorem);
(ii) SK-CAP structure of $\Phi_{10} = \Delta_5^2$
(Andrianov~$1974$; Arthur~$2013$ Thm.~$1.5.2$
\ClaimStatusTheorem);
(iii) metaplectic spin-character identification
$\nu_{\Delta_5}|_{\Gamma_{\mathrm{torsion}}} = \chi_{\mathrm{spin}}$
(Gritsenko--Nikulin~$1998$ \S$3$ \ClaimStatusTheorem).

The closure is \emph{reducible}: no genuine new theorem is needed;
the Hashimoto trace is forced by the three ingredients. The Wave-13~P5
wording ``trace $= 1$'' was a conflation with multiplicity-one
($\dim=1$); the correct trace is
eq.~\eqref{eq:wv14-k1-twisted-trace}, a transcendental quantity at
generic $p$ (since $a_p(\Delta_{E_6})$ is a Ramanujan-$\tau$-type
coefficient of weight $16$, $|a_p(\Delta_{E_6})| \le 2 p^{15/2}$ by
Deligne $1974$).

Arthur multiplicity input is present but non-essential: the
multiplicity formula Thm.~$1.5.2$ (Arthur~$2013$) delivers global
multiplicity one for the SK-CAP packet
$\psi_{\Phi_{10}} = \varphi_{\Delta_{E_6}} \boxtimes \mathrm{Sym}^1$,
and the metaplectic descent from $\Phi_{10}$ to $\Delta_5$ is the
double cover $\widetilde\Sp_4(\Q) \to \Sp_4(\Q)$ with
character $\nu_{\Delta_5}$ (Weissauer~$2009$ \S$3$). Ikeda-lift input
is unnecessary: $\Delta_5$ is \emph{not} an Ikeda lift (weight $5 < 2
\cdot 2$ threshold) but $\Delta_{10}$ fails to be an Ikeda lift too
(SK-CAP, not Ikeda $g=2$ image of a weight-$18$ $\mathrm{SL}_2$ form,
which would be genus $g=2$ from weight $2k-2 = 18$; direct check
fails, confirming SK-CAP-only).

\emph{Residual.} The closure of O1 upgrades Wave-13~P5
\ClaimStatusOpen\ to \ClaimStatusReducible, with the three reduction
ingredients all \ClaimStatusTheorem. The four $c=-214$ derivations of
Wave-11 are pinned to a \emph{common metaplectic scalar}: the
Gritsenko-convention $C_{\mathrm{Gr}} = 1$ on $a(T_0)$ and the Hecke
eigenvalue eq.~\eqref{eq:wv14-k1-twisted-trace} together force the
normalisation match across all four constructions, since all four
derivations must produce the same $\lambda_p(\Delta_5)$ up to the
finite sign $\chi_{\mathrm{spin}}(p)$, and the sign is pinned by the
Heegner-spin convention (Bruinier $2002$ Thm~$5.1$ regulator
$+1$ at $H_1$). The remaining P5 opens are O2 (full-$qq$-character
depth), O3 ($\mathcal T[A_1, \Sigma_{0,24}]$ Schur-index),
O4 (twisted-twined cocycle), O5 ($\hCS$-VOA realisation); none
obstructs O1.

\paragraph{Numerical check.} At $p = 2$: $a_2(\Delta_{E_6}) = 216$
(Ramanujan-type coefficient, weight $16$, standard table);
$\lambda_2(\Phi_{10}) = 216 + 256 + 512 = 984$;
$\lambda_2(\Delta_5) = \pm\sqrt{984}$, which is
\emph{not} rational. Since $\lambda_p(\Delta_5)$ must be an algebraic
integer in the Hecke eigenvalue field, this forces
$\lambda_p(\Delta_5) \in \Z[\sqrt{984}]$ or a quadratic extension
thereof; the Heegner-spin sign $\chi_{\mathrm{spin}}(2) = +1$ (Gritsenko
$1994$ Table) picks the positive root. This is consistent with the
Arthur-tempered Ramanujan bound $|\lambda_p(\Delta_5)| \le 4 p^{7/2}
= 4 \cdot 2^{7/2} \approx 45.3$, and $\sqrt{984} \approx 31.4 < 45.3$.

\emph{Cross-volume note.} Vol~I landscape-census: the chiral-side
$\kappa_{\mathrm{BKM}}(\Delta_5) = c_1(0)/2 = 5$ is unchanged (this is
the BKM weight of the Borcherds product, independent of the Hecke
eigenvalue). Vol~II cross-check: the $\mathsf{SC}^{\mathrm{ch,top}}$
realisation of $\mathbf H_{\Delta_5}$ on $K3 \times E$ reads
eq.~\eqref{eq:wv14-k1-twisted-trace} as the trace of Frobenius at $p$
on the $\ell$-adic realisation, with motivic weight $2k - 3 = 7$ and
spinor $L$-factor $(1 - \lambda_p(\Delta_5) p^{-s} + \cdots)^{-1}$.
The square-root structure matches Vol~III
\texttt{chapters/connections/geometric\_langlands.tex}
eq.~\eqref{eq:gl-gritsenko-additive}: no formula rewrite required
across volumes, the present note matches the existing identity
$\lambda_p^{\Delta_5} = \chi_{\mathrm{spin}}(p)\sqrt{\lambda_p(\Delta_{10})}$
already recorded in Vol~III.

\emph{Primary sources.} Hashimoto, \emph{Math.\ Ann.}~$262$ ($1983$)
$497$--$549$ (trace formula, twisted); Hashimoto, \emph{Proc.\ Symp.\
Pure Math.}~$86$ ($2012$), finite-conductor paramodular;
Ibukiyama, \emph{Comment.\ Math.\ Univ.\ St.\ Pauli}~$56$ ($2007$);
Ibukiyama--Wakatsuki, arXiv:$1305.4886$ ($2014$);
Gritsenko, \emph{St.\ Petersburg Math.\ J.}~$11$ ($1999$);
Gritsenko--Nikulin, arXiv:alg-geom/$9612002$ ($1997$) and
\emph{Amer.\ J.\ Math.}~$120$ ($1998$); Igusa, \emph{Amer.\ J.\
Math.}~$86$ ($1964$); Arthur, \emph{The Endoscopic Classification
of Representations} ($2013$) Thm.~$1.5.2$; Weissauer, \emph{Endoscopy
for $\mathrm{GSp}(4)$}, LNM~$1968$ ($2009$).
